\documentclass[12pt]{article}

% arXiv categories: cs.AI cs.AR q-bio.NC
\usepackage{amsmath, amssymb, amsthm}
\usepackage{geometry}
\usepackage{cite}
\usepackage{enumitem}
\usepackage{hyperref}  % load last

\geometry{margin=1in}

\newtheorem{definition}{Definition}
\newtheorem{proposition}{Proposition}
\newtheorem{conjecture}{Conjecture}
\newtheorem{remark}{Remark}

\title{Reconfigurable Dataflow Topology and the Hardware-Grounding\\
Alignment Problem:\\
A Structural Reformulation and Research Agenda\\
for CGLA-class Architectures}

\author{Minamo Minamoto\\
\small Independent\\
\small \texttt{minamominamoto4f5683f6@gmail.com}}

\date{April 9, 2026}

\begin{document}

\maketitle

\begin{abstract}
Current AI accelerator design is optimized for the computational
demands of large language models, but a structural question has
received insufficient attention: does the topology of a hardware
architecture constrain the class of sensory groundings it can
faithfully implement? We propose a reformulation of this question
using a topological framework introduced in a companion paper
\cite{minamoto2026topology}, which formalizes symbol grounding as a
structural property of sensory space and introduces the
Gromov--Hausdorff cognitive distance $d_{\mathrm{cog}}$ as a
measure of grounding alignment between cognitive systems.

Within this framework, we argue that fixed-dataflow accelerators
---including TPUs, GPUs, and memory-integrated SoCs such as the
Apple M-series---impose a topological constraint on the class of
sensory groundings they can implement, incurring a strictly
positive $d_{\mathrm{cog}}$ lower bound for sensory spaces not
embeddable in their fixed computational graph topology. We further
propose that CGLA-class reconfigurable architectures are
\emph{candidates} for the architectural property required for
universal sensory grounding, though this claim rests conditionally
on an unproven conjecture (Section~4.2) whose verification or
refutation is the central open obligation of this research program.

This paper does not present a completed theory. It presents a
structural reformulation of the hardware-grounding alignment
problem, identifies the key unresolved challenge for CGLA-class
architectures, and proposes a research agenda connecting hardware
topology to sensory topology. The necessary empirical and formal
completion of this program---including verification of CGLA
topological universality and empirical measurement of
$d_{\mathrm{cog}}$ across hardware substrates---requires
collaboration between cognitive theory and hardware architecture
research. We close with five open problems at this intersection.
\end{abstract}

\tableofcontents

\newpage

%-------------------------------------------------------------------
\section{Introduction}
%-------------------------------------------------------------------

The dominant AI accelerator paradigm is shaped by the computational demands
of transformer-based large language models: massive matrix multiplication,
high memory bandwidth, and fixed dataflow optimized for dense linear algebra.
Google's Tensor Processing Unit (TPU) represents the extreme of this design
philosophy---a systolic array architecture in which the dataflow between
processing elements is fixed at design time \cite{jouppi2017}. This fixity
yields exceptional performance per watt for the specific computational
graphs of transformer inference, but it purchases efficiency at the cost of
architectural flexibility.

A parallel line of inquiry has asked what \emph{cognitive} requirements
AI hardware should satisfy. If the goal of an AI system is not merely to
process text efficiently, but to implement something resembling genuine
understanding---grounded in sensorimotor experience and capable of concept
formation across diverse domains---then the hardware substrate may matter
in ways that current accelerator design does not address.

In a companion paper \cite{minamoto2026topology}, we proposed a structural
reformulation of symbol grounding: the topology of the sensory space
in which grounding occurs may determine the geometry of the resulting
cognitive representations, and two agents' grounding-induced topologies
being homeomorphic is proposed as a sufficient condition for full
intensional concept sharing. This reformulation of Harnad's symbol
grounding problem \cite{harnad1990} from a binary to a structural
condition, if correct, has direct implications for hardware design.

The central proposal of the present paper is:

\begin{quote}
\emph{Fixed-dataflow architectures impose a topological restriction on
implementable sensory spaces---this is the paper's primary formal
claim (Proposition~1, conditional on Definition~5). We further
propose that CGLA-class reconfigurable architectures are strong
candidates for the architectural property required for universal
sensory grounding, but this claim rests on Conjecture~1, which
remains unproven. This paper's goal is to make these structural
claims precise, identify what verification would require, and
propose a research agenda for CGLA-class architecture as a
candidate substrate for grounding-faithful AI systems.}
\end{quote}

We develop this proposal in three steps. Section~2 reviews the
framework from \cite{minamoto2026topology}. Section~3 argues that fixed
dataflow implies topological restriction under Definition~5.
Section~4 presents CGLA-class reconfigurability as a structural
candidate and compares it qualitatively to FPGAs and fixed-dataflow
alternatives. Section~5 discusses implications and introduces
$d_{\mathrm{cog}}$ as a hardware selection criterion. Section~6
lists open problems that the CGLA research community is best
positioned to resolve.

\begin{remark}[On the status of this paper]
This paper is a \emph{foundational position paper}, not a completed
theory. Its contribution is conceptual: it proposes a structural
reformulation of the hardware-grounding alignment problem, identifies
the topological restriction imposed by fixed-dataflow architectures,
and specifies what CGLA-class architectures would need to demonstrate
in order to satisfy the universality condition. The formal completion
of the program---proof or disproof of Conjecture~1, empirical
measurement of $d_{\mathrm{cog}}$, and verification that CGLA
configurations faithfully implement biological sensory
topologies---requires collaboration between cognitive theory and
hardware architecture research, and is beyond the scope of this paper
alone.
\end{remark}

%-------------------------------------------------------------------
\section{Background: Sensory Space Topology and Grounding}
%-------------------------------------------------------------------

We briefly summarize the framework of \cite{minamoto2026topology}, to which we
refer the reader for full definitions and proofs.

\begin{definition}[Sensory Space \cite{minamoto2026topology}]
The \emph{sensory space} of a cognitive agent $\mathcal{A}$ equipped with
transduction channels $\{T_1, \ldots, T_n\}$ is the product space
\[
\mathcal{S}(\mathcal{A}) = \prod_{i=1}^{n} A_i
\]
equipped with the product topology, where each $A_i$ carries the topology
induced by the neural activation metric of channel $T_i$.
\end{definition}

\begin{definition}[Grounding Map \cite{minamoto2026topology}]
Let $\mathcal{K}(\mathcal{S}(\mathcal{A}))$ denote the hyperspace of
nonempty closed subsets of $\mathcal{S}(\mathcal{A})$, equipped with the
Vietoris topology \cite{vietoris1922}. A \emph{grounding map} is a
function
\[
G: \Sigma \to \mathcal{K}(\mathcal{S}(\mathcal{A}))
\]
where $\Sigma$ is a symbol set and $G(\sigma)$ is the closed set of
sensory states associated with symbol $\sigma$.
\end{definition}

\begin{definition}[Grounding-Induced Topology \cite{minamoto2026topology}]
The \emph{grounding-induced topology} $\tau_{\mathcal{A}}$ on $\Sigma$
is the initial topology with respect to $G$: the coarsest topology making
$\sigma \mapsto G(\sigma)$ continuous with respect to the Vietoris
topology on $\mathcal{K}(\mathcal{S}(\mathcal{A}))$.
\end{definition}

The key result of \cite{minamoto2026topology} is that homeomorphism of
grounding-induced topologies is \emph{sufficient} for complete intensional
concept sharing (Conjecture~1 of \cite{minamoto2026topology}), while topological
inequivalence implies incommensurability of intensional content. The
Gromov--Hausdorff cognitive distance
\[
d_{\mathrm{cog}}(\mathcal{A}, \mathcal{B})
= d_{GH}\!\left((\Sigma, \tau_{\mathcal{A}}),\, (\Sigma, \tau_{\mathcal{B}})\right)
\]
measures the degree of grounding alignment between two agents, with
$d_{\mathrm{cog}} = 0$ corresponding to the homeomorphic (complete sharing)
case \cite{gromov1981}.

%-------------------------------------------------------------------
\section{Fixed Dataflow as Topological Restriction}
%-------------------------------------------------------------------

\subsection{Computational Graphs and Their Topology}

A neural network computation can be represented as a directed acyclic graph
$\mathcal{G} = (V, E)$ where vertices $V$ are processing elements and
directed edges $E$ represent dataflow. The \emph{computational graph
topology} is the topology induced on $V$ by the adjacency structure of
$\mathcal{G}$.

\begin{definition}[Fixed Dataflow Architecture]
An architecture is \emph{fixed-dataflow} if the computational graph
$\mathcal{G}$ is determined at design time and cannot be reconfigured
during operation. The topology of $\mathcal{G}$ is therefore a constant
of the hardware.
\end{definition}

TPUs are fixed-dataflow in this sense: the systolic array defines a rigid
grid topology in which data flows in a fixed pattern optimized for
matrix multiplication. Alternative operations---including those arising
from non-matrix sensory modalities---must be mapped onto this fixed
topology, potentially at significant computational cost or loss of
structural fidelity.

\subsection{The Topological Restriction Proposition}

The connection between sensory space topology and computational graph
topology is the following. Any implementation of a grounding map $G$ on
hardware requires a mapping from the topological structure of
$\mathcal{S}(\mathcal{A})$ to the topological structure of the
computational graph $\mathcal{G}$. We formalize this as follows.

\begin{definition}[Topological Implementation Map]
A \emph{topological implementation map} from sensory space
$(\mathcal{S}(\mathcal{A}), \tau_{\mathcal{A}})$ to a hardware
architecture $\mathcal{H}$ with computational graph
$(\mathcal{G}, \tau_{\mathcal{G}})$ is a continuous injection
\[
\phi: (\Sigma, \tau_{\mathcal{A}}) \hookrightarrow
      (\mathcal{G}, \tau_{\mathcal{G}}).
\]
The implementation is \emph{faithful} if $\phi$ is a homeomorphism
onto its image, i.e., both $\phi$ and $\phi^{-1}$ preserve the
topological structure of the grounding-induced symbol space.
\end{definition}

\begin{remark}[Conceptual justification of Definition~5]
Definition~5 maps between two spaces of fundamentally different
character: $(\Sigma, \tau_{\mathcal{A}})$ is a semantic space---the
symbol set equipped with the topology induced by sensorimotor grounding
---while $(\mathcal{G}, \tau_{\mathcal{G}})$ is a physical space---the
processing element graph of a hardware chip. This cross-domain mapping
requires justification.

The conceptual bridge is the following: any physical implementation
of a grounding map $G: \Sigma \to \mathcal{K}(\mathcal{S}(\mathcal{A}))$
must assign each symbol $\sigma \in \Sigma$ to a set of computational
states that represent $G(\sigma)$, and must implement the relational
structure among symbols as relational structure among computational
processes. Symbols that are topologically close in $\tau_{\mathcal{A}}$
---meaning their grounding sets are similar in sensory space---should
correspond to computationally related processes in $\mathcal{G}$.
If the computational graph lacks the relational structure to place such
processes in proximity, the topological relationships among grounded
symbols are distorted. Definition~5 formalizes this requirement as
continuity: $\phi$ preserves the topological relationships of the
semantic space in the computational graph.

The faithfulness condition (homeomorphism onto image) further requires
that no relational structure is lost in either direction: neither is
$\tau_{\mathcal{A}}$-proximity collapsed into $\tau_{\mathcal{G}}$-
distance (loss of discrimination), nor is $\tau_{\mathcal{G}}$-proximity
imposed on $\tau_{\mathcal{A}}$-distant symbols (spurious conflation).

We acknowledge that this definition conflates semantic topology and
physical topology in a way that requires empirical grounding---namely,
the specification of what metric on $\mathcal{G}$ is relevant for
cognitive purposes. This is OP1; the formal adequacy of Definition~5
is therefore conditional on its resolution.
\end{remark}

The following restriction holds as a direct consequence of
Definition~5; its force depends on the adequacy of that definition,
which we identify as the most fundamental open problem of this
framework (OP1).

\begin{proposition}[Topological Restriction of Fixed Dataflow
  (under Definition~5)]
Let $\mathcal{H}$ be a fixed-dataflow architecture with computational
graph topology $\tau_{\mathcal{G}}$. Then $\mathcal{H}$ can implement
only grounding maps $G$ whose induced topology $\tau_{\mathcal{A}}$ is
homeomorphic to a subspace of $(\mathcal{G}, \tau_{\mathcal{G}})$.
Sensory spaces whose topology is not homeomorphic to any subspace of
$(\mathcal{G}, \tau_{\mathcal{G}})$ cannot be faithfully grounded on
$\mathcal{H}$.
\end{proposition}

\begin{remark}
We emphasize the distinction between \emph{computational universality}
and \emph{topological faithfulness}. A fixed-dataflow architecture is
Turing-complete in the sense that it can compute any computable
function; the restriction here is not on computability but on structural
preservation. When a grounding map $G$ whose sensory topology
$\tau_{\mathcal{A}}$ is not embeddable in $\tau_{\mathcal{G}}$ is
executed on $\mathcal{H}$, the computation proceeds correctly but the
topological relationships among grounded symbols are distorted.
Zero-padding and format-conversion workarounds preserve the numerical
output of individual computations while destroying the relational
structure of the grounding space. This distortion is therefore
structural, not merely a performance overhead: it implies
$d_{\mathrm{cog}} > 0$ as a \emph{lower bound} on the achievable
cognitive distance, not a recoverable inefficiency.
\end{remark}

\subsection{Concrete Case: TPU and Chemical Sensory Space}

Consider the sensory space of \textit{Caenorhabditis elegans}: a
low-dimensional chemical topology dominated by volatile compound gradients
\cite{bargmann2006}. The natural computational graph for implementing
chemotaxis-based grounding is a sparse, gradient-following structure with
few processing elements and sequential update dynamics \cite{white1986}.

The TPU's systolic array, by contrast, is optimized for dense $n \times n$
matrix products. Mapping a sparse gradient-following computation onto a
dense matrix architecture requires either zero-padding (wasting compute)
or restructuring the computation to fit the matrix template (distorting
the grounding topology). Neither preserves the topological structure of
the original sensory space.

This is a concrete instance of topological restriction: the TPU's fixed
topology is well-matched to transformer attention but poorly matched to
chemical gradient computation. Any agent implemented on a TPU that
attempts to ground concepts in a chemical sensory space will suffer
elevated $d_{\mathrm{cog}}$ relative to a native chemical-topology
implementation.

\subsection{Contrast Case: GPU and Kernel-Level Pseudo-Flexibility}

Graphics Processing Units (GPUs) appear more flexible than TPUs:
CUDA kernels can be selected at runtime to implement diverse
computational patterns, and the programmer can target different
memory access patterns and parallelization strategies depending on
the workload. This might suggest that GPUs occupy an intermediate
position in topological flexibility between TPUs and CGLA-class
architectures.

We argue this appearance is misleading. GPU flexibility resides
at the \emph{software layer}---kernel selection determines which
computation executes on the hardware, but does not change the
hardware topology itself. The underlying physical architecture is
a fixed grid of Streaming Multiprocessors (SMs) connected through
a hierarchical memory system (L1 cache, L2 cache, HBM). This
hardware topology is determined at design time and is not
reconfigurable during operation.

Under Definition~3, GPUs are therefore fixed-dataflow in the
relevant sense: the computational graph topology $\tau_{\mathcal{G}}$
of the hardware remains constant across kernel invocations.
A sparse gradient-following computation (such as the
\textit{C. elegans} chemotaxis example of Section~3.2) mapped
onto a GPU requires the same structural translation as on a TPU:
the sparse graph must be padded or restructured to fit the
dense SM grid topology. The resulting $d_{\mathrm{cog}}$ lower
bound is therefore qualitatively identical to the TPU case,
differing in magnitude depending on the degree of sparsity
mismatch rather than in kind.

This distinction---between software-layer flexibility and
hardware-topology reconfigurability---is the critical boundary
separating fixed-dataflow and reconfigurable-dataflow
architectures for the purposes of grounding fidelity. The CGLA's
reconfigurable interconnect changes the actual computational
graph topology $\tau_{\mathcal{G}}$; GPU kernel selection does not.

We note that Field-Programmable Gate Arrays (FPGAs) also offer
hardware-level reconfigurability and thus do not fall under the
fixed-dataflow restriction of Proposition~1. The comparative analysis
of $d_{\mathrm{cog}}$-minimization between CGLA-class and FPGA
architectures is deferred to OP5; the present argument is that
reconfigurability is a necessary architectural property for universal
grounding, not that CGLA is its unique instantiation.

\subsection{Extended Case: Apple M-series and Memory-Integrated Fixed Dataflow}

Apple's M-series System-on-Chip (SoC) architecture represents a
distinct and instructive case that differs from both TPU and GPU in
its integration strategy. The M5 chip \cite{apple2025m5} integrates
a 16-core Neural Engine, a 10-core GPU with a dedicated Neural
Accelerator in each core, and a CPU, all sharing a unified memory
pool via a high-bandwidth interconnect (153 GB/s). This unified
memory architecture eliminates data transfer overhead between
processing units---a genuine engineering advantage over discrete
chip configurations.

However, the topological analysis of the present framework reveals
that memory integration and dataflow reconfigurability are orthogonal
properties. The M-series Neural Engine is a fixed-dataflow
accelerator: its computational graph topology is determined at
design time and optimized for matrix-based neural network operations.
The Neural Accelerators within each GPU core are similarly
fixed-topology units. The unified memory architecture addresses the
\emph{bandwidth} bottleneck between processing units, but does not
address the \emph{topological mismatch} between the hardware's fixed
computational graph structure and the target sensory space topology.

Concretely: an M5 implementation of \textit{C.~elegans} chemotaxis
faces the same topological restriction as a TPU implementation.
The sparse, gradient-following computational graph required by the
chemotaxis loop must be mapped onto fixed matrix-operation units,
incurring the same structural distortion that elevates
$d_{\mathrm{cog}}$. The unified memory reduces the cost of
transferring intermediate representations between CPU, GPU, and
Neural Engine, but the representations themselves are computed in
fixed-topology units that impose the same homeomorphism constraint
as Proposition~1.

\begin{remark}[Memory Integration vs.\ Topological Flexibility]
This distinction is architecturally fundamental. Memory-integrated
SoC designs (Apple M-series, Qualcomm Snapdragon) optimize the
\emph{communication} topology between fixed-dataflow processing
units. CGLA-class reconfigurable architectures optimize the
\emph{computational} topology of the processing units themselves.
For the purposes of $d_{\mathrm{cog}}$-minimization, only the
latter is relevant: the cognitive distance between two grounding
implementations is determined by the computational graph topology
available to implement the grounding map $G$, not by the
memory bandwidth available to transfer activations between stages.
A system with high memory bandwidth and fixed computational topology
will achieve lower $d_{\mathrm{cog}}$ than a low-bandwidth system
with the same fixed topology---but a system with reconfigurable
computational topology can achieve $d_{\mathrm{cog}} \to 0$ for
sensory spaces within its configuration space, regardless of
memory bandwidth.
\end{remark}

%-------------------------------------------------------------------
\section{CGLA-class Reconfigurability as Topological Flexibility}
%-------------------------------------------------------------------

\subsection{CGLA Architecture}

The Coarse-Grained Linear Array (CGLA) is a reconfigurable computing
architecture in which the dataflow between processing elements can be
modified at runtime \cite{akabe2025, ando2026}. Unlike the systolic
array's fixed grid, CGLA allows the computational graph topology
$\tau_{\mathcal{G}}$ to be reconfigured to match the structure of the
workload.

Recent work has demonstrated CGLA implementations of LLM inference
\cite{ando2026_llm}, attention acceleration \cite{takeuchi2025}, and
energy-efficient whisper ASR \cite{ando2025_whisper}, showing that the
architecture can handle diverse computational structures within a single
hardware substrate. The key property for our purposes is not performance
per se, but \emph{topological flexibility}: the ability to present
different computational graph topologies to different workloads.

We note that other reconfigurable and neuromorphic architectures---
including Intel's Loihi~2 \cite{davies2021}, SpiNNaker
\cite{furber2014}, and BrainScaleS \cite{schemmel2010}---also offer
topological flexibility via configurable spike-routing networks. The
present argument applies to the class of reconfigurable-dataflow
architectures broadly; CGLA is the primary instantiation we consider
because its hardware-efficiency properties under diverse computational
workloads have been most recently documented \cite{akabe2025,
ando2025_whisper}. A systematic comparative analysis of
$d_{\mathrm{cog}}$-minimization across these architectures is
deferred to OP5.

\subsection{Topological Flexibility and Universal Grounding}

\begin{definition}[Topologically Universal Architecture]
An architecture $\mathcal{H}$ is \emph{topologically universal} with
respect to a class $\mathcal{C}$ of sensory spaces if for every
$\mathcal{S} \in \mathcal{C}$, there exists a configuration of
$\mathcal{H}$ whose computational graph topology is homeomorphic to the
grounding-induced topology of $\mathcal{S}$.
\end{definition}

\begin{conjecture}[CGLA Topological Universality]
CGLA-class architectures are topologically universal with respect to the
class of sensory spaces arising from biological transduction channels
(chemical, visual, mechanoreceptive, proprioceptive, and their
cross-modal products).
\end{conjecture}

\begin{remark}[On the strength of the homeomorphism condition]
The universality conjecture requires exact homeomorphism between
CGLA computational graph topologies and biological sensory topologies.
This is a strong condition that may not hold in practice, for two
reasons.

First, biological sensory spaces are continuous, high-dimensional,
and fiber-bundle-structured in many organisms (Section~3.2 of
\cite{minamoto2026topology}), while CGLA computational graphs are finite and
discrete. A faithful homeomorphism in the strict sense therefore
requires discretization, which necessarily involves approximation.

Second, even if CGLA's configuration space is dense in the metric
space of sensory topologies (OP5), there may be sensory topologies
that require configurations exceeding CGLA's physical connectivity
constraints.

We therefore propose a weakened version of the universality condition
as an alternative target for verification: rather than exact
homeomorphism, require that CGLA can achieve $d_{\mathrm{cog}}$
below a threshold $\epsilon > 0$ for all sensory spaces in the
biological class. This $\epsilon$-universality condition is more
empirically tractable and biologically meaningful, since perfect
grounding fidelity is not required for functional cognitive
equivalence. We flag this weakening as part of OP2.
\end{remark}

The conjecture is supported by the following structural argument. The
CGLA's reconfigurable interconnect allows the processing element graph
to be set to any sparse directed graph within the physical connectivity
constraints of the chip. The sensory spaces of biological organisms, while
high-dimensional, are structured by the specific connectivity patterns of
neural circuits---which are themselves sparse and locally connected
\cite{white1986, scheffer2020}. A reconfigurable architecture that can
present sparse, locally-connected computational graphs is therefore
structurally aligned with biological sensory topologies in a way that
dense matrix architectures are not.

\begin{remark}[Conditional Necessity]
The paper's central claim---reconfigurable dataflow topology as a
\emph{necessary} condition---rests conditionally on this conjecture.
What follows from Proposition~1 alone is that fixed-dataflow
architectures impose a strictly positive lower bound on
$d_{\mathrm{cog}}$ for sensory spaces not embeddable in
$\tau_{\mathcal{G}}$. The necessity claim then holds as follows: if
Conjecture~1 holds, then CGLA-class reconfigurability achieves a lower
bound of zero, making it the unique class of physical architectures
capable of faithful universal grounding. Proof of Conjecture~1 (OP2)
is therefore the most pressing remaining obligation of this research
program.
\end{remark}

\subsection{CGLA vs.\ FPGA: A Qualitative Comparison}

As noted in Section~3.3, FPGAs also offer hardware-level
reconfigurability and are therefore not subject to the fixed-dataflow
restriction of Proposition~1. This raises a legitimate question: is
CGLA's advantage over fixed-dataflow architectures shared equally by
FPGAs, and if so, what is CGLA's distinctive contribution?

We offer a qualitative comparison on three dimensions relevant to
$d_{\mathrm{cog}}$-minimization.

\textbf{Reconfiguration granularity.} FPGAs achieve reconfigurability
via lookup tables (LUTs) and programmable interconnects at the
gate level, enabling arbitrary logic functions. CGLA achieves
reconfigurability at the level of coarse-grained functional units
(arithmetic and logical operations) connected by a reconfigurable
routing network \cite{akabe2025}. For the purpose of implementing
sensory grounding topologies---which are structured by sparse,
directed graphs of moderate complexity rather than arbitrary Boolean
circuits---CGLA's coarse-grained reconfigurability is well-matched
to the target structure. FPGA's fine-grained LUT-level programming
introduces significant overhead (both area and latency) for the
kind of structured sparse computations that biological sensory
topologies require. In this sense, CGLA is a ``right-sized''
reconfigurable architecture for grounding-topology implementation,
whereas FPGA is over-specified for this task and incurs avoidable
energy costs.

\textbf{Energy efficiency.} CGLA's empirically demonstrated energy
efficiency \cite{akabe2025, ando2025_whisper} reflects this
granularity match: reconfiguring at the functional-unit level
consumes less power than reconfiguring at the gate level for
computations that do not require gate-level flexibility. For
grounding implementations intended to run continuously in embodied
systems---where energy efficiency is a hard constraint---this
distinction is practically significant.

\textbf{Scope of the present argument.} The claim of the present
paper is that \emph{reconfigurable dataflow topology is necessary}
for universal grounding, not that CGLA is the unique architecture
satisfying this condition. FPGAs satisfy the topological
reconfigurability condition and should in principle achieve
$d_{\mathrm{cog}} \to 0$ for sensory spaces within their
configuration space. The comparative analysis of
$d_{\mathrm{cog}}$-minimization efficiency between CGLA and FPGA
implementations is deferred to OP5; the present argument establishes
the necessary property, and CGLA is the primary instantiation we
consider because its efficiency under diverse workloads has been most
recently and rigorously documented \cite{akabe2025, ando2025_whisper}.

%-------------------------------------------------------------------
\section{Implications}
%-------------------------------------------------------------------

\subsection{Hardware Selection as Cognitive Design}

The framework developed here implies that hardware selection is not merely
a performance engineering decision but a cognitive design decision. An AI
system intended to implement human-like grounding requires hardware whose
computational graph topology can approximate the topology of human sensory
space. A system grounded in a chemical sensory space (like a robotic nose)
requires hardware whose topology approximates a sparse gradient-following
graph.

The cognitive distance $d_{\mathrm{cog}}$ provides a principled criterion
for hardware selection: given a target sensory topology $\tau_{\mathcal{A}}$,
select the hardware architecture $\mathcal{H}$ that minimizes
$d_{\mathrm{cog}}$ between $\tau_{\mathcal{A}}$ and the best achievable
computational graph topology of $\mathcal{H}$.

For transformer-based LLMs with attention-pattern topology, this criterion
favors TPUs. For systems intended to implement diverse sensory groundings,
this criterion favors CGLA-class architectures.

\subsection{Neuro-Symbolic Alignment}

Neuro-symbolic AI systems combine a neural perception layer with a
symbolic reasoning layer \cite{garcez2022}. The topological alignment
condition of \cite{minamoto2026topology}---that the symbolic layer's induced
topology should be consistent with the neural layer's grounding
topology---extends naturally to hardware: the hardware topology of the
neural layer should be consistent with the sensory topology, and the
hardware topology of the symbolic layer should be consistent with the
relational structure of the symbol set.

CGLA-class architectures, capable of presenting different topologies to
different layers of a neuro-symbolic system, are therefore structurally
suited to implementing topological alignment at the hardware level.

\subsection{Energy Efficiency as a Consequence of Topological Matching}

A secondary implication concerns energy efficiency. When hardware topology
is mismatched to computational topology, energy is expended on structural
translation: zero-padding, reshaping, and format conversion operations
that carry no semantic content. Topological matching eliminates this
overhead.

This connects to CGLA's empirically demonstrated energy efficiency
advantages \cite{akabe2025}: the efficiency gain is not merely an
engineering achievement but a structural consequence of reduced topological
mismatch between the computation and the hardware.

%-------------------------------------------------------------------
\section{Pilot Evidence: RSA as $d_\mathrm{cog}$ Proxy in LLMs}
\label{sec:pilot}
%-------------------------------------------------------------------

\begin{remark}[Which component of $d_\mathrm{cog}$ does RSA measure?]
Under the persistence signature formulation of
\cite[Def.~\ref{def:dcog-signature}]{minamoto2026topology}, 
$d_\mathrm{cog}$ is a tuple whose first two components 
($d_\mathrm{GH}$ and $d_\mathrm{bn}^{(0)}$) are computable from an 
RSA-derived metric on $\Sigma$. The second-order Spearman 
$\rho_\mathrm{cog}$ used below is a monotone proxy for 
$d_\mathrm{bn}^{(0)}$ --- it tracks merge-structure alignment, not 
loop structure. Higher-order components $d_\mathrm{bn}^{(k)}$ 
($k \geq 1$) are not resolved at the twenty-nine--concept inventory 
used here (the detection threshold is $N \gtrsim 24$; see 
\cite[App.~\ref{app:dcog-h1}]{minamoto2026topology}). The pilot 
below therefore substantiates the $H_0$-level proxy claim; the 
$H_1$-level remains open.
\end{remark}

OP4 proposes RSA as a tractable proxy for $d_\mathrm{cog}$ in toy
hardware experiments.
We report a pilot study applying this proxy to language model
training checkpoints, providing a software-level proof of concept
for the RSA-based $d_\mathrm{cog}$ measurement protocol.

\paragraph{Setup.}
We computed RSA dissimilarity matrices over $n = 29$ probe concepts
at hidden layers $\ell \in \{2, 4, 6\}$ of Pythia-70m
\cite{biderman2023} at three training checkpoints: step~0 (random
initialization), step~1000 (early training), step~143000 (fully
trained reference).
Three prompt templates were averaged to reduce surface-form bias.
Structural alignment to the reference was measured as the
second-order Spearman correlation $\rho_\mathrm{cog}$, which OP4
proposes as the operative proxy for $d_\mathrm{cog}$.

\paragraph{Results.}
$\rho_\mathrm{cog}$ increases monotonically with training at all
layers (L6: step-0 $= 0.017$; step-1000 $= 0.535$), confirming
that RSA captures a meaningful structural signal that evolves
predictably with training.
Cross-model replication with GPT-2 \cite{radford2019language}
shows $\rho_\mathrm{cog}(\text{GPT-2}, \text{Pythia-trained})
\approx 0.42$ versus $\rho_\mathrm{cog}(\text{GPT-2},
\text{Pythia-step-0}) \approx 0.10$, demonstrating that
$\rho_\mathrm{cog}$ discriminates structurally distinct
representations.

\paragraph{Relevance to hardware experiments.}
This result supports the validity of the RSA proxy for the
hardware $d_\mathrm{cog}$ measurement proposed in OP3--OP4.
Specifically, it demonstrates that:
(i) $\rho_\mathrm{cog}$ is sensitive to structural differences
in representational geometry at fixed architecture;
(ii) the proxy is stable across prompt formulations when averaged;
(iii) cross-architecture comparison (GPT-2 vs. Pythia) is
feasible with the same protocol.

The analogous hardware experiment---comparing RSA dissimilarity
matrices of grounded representations on TPU, GPU, and CGLA
substrates---would use the same $\rho_\mathrm{cog}$ protocol,
substituting hardware-level activation traces for model
checkpoint activations.
The prediction of Section~5 is that CGLA implementations yield
higher $\rho_\mathrm{cog}$ to the target topology than fixed-dataflow
substrates at matched computational budget.
All code is available at
\url{https://github.com/minamominamoto/grounding-without-reality}.

%-------------------------------------------------------------------
\section{Open Problems}
%-------------------------------------------------------------------

\begin{enumerate}[label=\textbf{OP\arabic*.}, leftmargin=*, itemsep=6pt]

\item \textbf{Formal definition of topological implementation map.}
    Define precisely the mapping from sensory space topology to
    computational graph topology. This requires specifying which
    properties of $\tau_{\mathcal{A}}$ are preserved under
    discretization onto a finite processing element graph, and which
    are necessarily lost.

\item \textbf{Proof or disproof of CGLA Topological Universality.}
    The conjecture in Section~4.2 requires a formal proof that
    CGLA's reconfigurable interconnect can present computational graph
    topologies homeomorphic to the grounding-induced topologies of
    biological sensory spaces. A disproof would require exhibiting a
    biologically motivated sensory topology that cannot be approximated
    by any CGLA configuration within given physical constraints.

\item \textbf{Empirical measurement of $d_{\mathrm{cog}}$ on hardware.}
    Given a fixed sensory grounding task (e.g., chemical gradient
    following in a simulated \textit{C. elegans} environment), measure
    $d_{\mathrm{cog}}$ between implementations on TPU, GPU, and
    CGLA-class hardware. The prediction of the present paper is that
    CGLA achieves the smallest $d_{\mathrm{cog}}$ at comparable
    energy expenditure.

\item \textbf{Toy model of $d_{\mathrm{cog}}$.}
    As a minimal concrete instantiation prior to full hardware
    experiments, consider a two-dimensional chemical sensory space
    $\mathcal{S}_{\mathrm{chem}}$ and a three-dimensional visual sensory
    space $\mathcal{S}_{\mathrm{vis}}$, each with synthetically generated
    grounding maps. Representational Similarity Analysis (RSA)
    \cite{kriegeskorte2008} provides an operational proxy for
    $d_{\mathrm{cog}}$: compute the dissimilarity matrix of grounded
    representations on each hardware substrate and compare via Spearman
    correlation to the target topology's dissimilarity matrix.

    \textbf{On the validity of RSA as a proxy for $d_{\mathrm{cog}}$.}
    RSA operates on pairwise dissimilarity matrices and is therefore
    sensitive to the ordinal structure of representational similarity
    rather than to absolute distances. Its Spearman correlation with a
    target dissimilarity matrix approximates $d_{\mathrm{cog}}$ under
    the assumption that the sensory space is \emph{locally isometric}
    in the regions of interest---that is, that the ratio of pairwise
    distances is preserved under the grounding map $G$ in a
    neighbourhood around each grounded concept. When this assumption
    holds, RSA similarity correlates inversely with $d_{\mathrm{cog}}$:
    a higher Spearman correlation indicates smaller topological
    distortion. When the sensory space is highly non-isometric (e.g.,
    fiber bundle structures as in \textit{Drosophila}), RSA may
    underestimate $d_{\mathrm{cog}}$ because it does not capture
    global topological invariants such as connectivity or homotopy
    type. In such cases, persistent homology methods would provide
    a more faithful proxy. For the purposes of the toy model (chemical
    and visual sensory spaces with product topology), the local
    isometry assumption is plausible and RSA is a tractable first
    approximation.

    The prediction is that, for non-matrix sensory spaces at matched
    computational budgets: (i) CGLA implementations yield higher RSA
    similarity to the target topology than both TPU and GPU implementations;
    (ii) TPU and GPU implementations yield qualitatively similar RSA
    similarity to each other, differing in magnitude but not in kind,
    since both are fixed at the hardware topology layer (Section~3.3);
    and (iii) the GPU--CGLA gap exceeds the TPU--GPU gap, because GPU
    software-layer flexibility does not reduce the structural translation
    cost that drives $d_{\mathrm{cog}}$.

\item \textbf{Extension to cross-modal sensory spaces.}
    The \textit{Drosophila} central complex exhibits cross-modal
    convergence that cannot be captured by a product topology
    \cite{scheffer2020}. Extending the framework to fiber bundle
    sensory spaces and determining whether CGLA's interconnect
    can represent fiber bundle computational graphs is an open problem.

\item \textbf{Hardware topology metric.}
    Define a metric on computational graph topologies analogous to
    $d_{\mathrm{cog}}$, and establish whether the CGLA's configuration
    space is dense in this metric space (i.e., whether it can
    approximate any target computational graph topology to arbitrary
    precision within physical constraints).

\end{enumerate}

%-------------------------------------------------------------------
\section{Conclusion}
%-------------------------------------------------------------------

We have proposed a structural reformulation of the hardware-grounding
alignment problem. Within the topological framework of
\cite{minamoto2026topology}, fixed-dataflow architectures incur a strictly
positive $d_{\mathrm{cog}}$ lower bound for sensory spaces not
embeddable in their fixed computational graph topology---this is the
paper's primary formal claim, conditional on Definition~5. We have
further proposed that CGLA-class reconfigurable architectures are
strong candidates for the architectural property required for
universal sensory grounding, though this candidacy rests on
Conjecture~1, which remains unproven.

This paper does not claim to have solved the hardware-grounding
alignment problem. It claims to have made the problem formally
precise. The research agenda that follows from this reformulation
has three components.

\textbf{For cognitive theory}: Provide a formal proof or disproof
of Conjecture~1 (OP2). Determine whether exact homeomorphism is the
correct universality condition, or whether the $\epsilon$-universality
weakening is more appropriate (Section~4.2). Formalize Definition~5
(OP1) so that the topological implementation map has a rigorous
operational meaning.

\textbf{For CGLA architecture research}: Characterize the
configuration space of CGLA's reconfigurable interconnect formally,
and determine whether it is dense in the metric space of biological
sensory topologies (OP5). Provide at least one concrete instantiation
of Definition~5---a specific CGLA configuration that implements the
\textit{C.~elegans} chemosensory topology---as a proof of concept.
Measure $d_{\mathrm{cog}}$ between CGLA and fixed-dataflow
implementations of the same sensory grounding task using RSA as
a proxy (OP3, OP4).

\textbf{For the intersection}: As AI systems are increasingly
deployed in embodied, multi-modal, and neuromorphic contexts, the
alignment between hardware topology and sensory topology will become
a principled design criterion. The cognitive distance $d_{\mathrm{cog}}$
provides the formal foundation for this criterion. Realizing it as
an engineering tool requires collaboration between the communities
that this paper attempts to bridge.

%-------------------------------------------------------------------
\section*{Acknowledgments}
%-------------------------------------------------------------------

No funding was received for this work. The author thanks the CGLA research
group at Nara Institute of Science and Technology for the foundational
hardware work cited here.

%-------------------------------------------------------------------
\begin{thebibliography}{99}

\bibitem{apple2025m5}
Apple Inc. (2025). Apple M5 chip. \textit{Wikipedia},
retrieved April 2026.
\texttt{https://en.wikipedia.org/wiki/Apple\_M5}
Note: No peer-reviewed publication describing M5 architecture
is available at time of writing; the Wikipedia article
aggregates specifications from Apple's technical documentation.
For the M-series unified memory architecture
as independently characterized, see \cite{brinkmann2025}.

\bibitem{brinkmann2025}
Brinkmann, M., et al. (2025). Apple vs.\ Oranges: Evaluating
the Apple Silicon M-Series SoCs for HPC Performance and
Efficiency. \textit{arXiv preprint arXiv:2502.05317}.

\bibitem{akabe2025}
Akabe, T., Le, V. T. D., \& Nakashima, Y. (2025). IMAX: A power-efficient
multilevel pipelined CGLA and applications. \textit{IEEE Access},
10.1109/ACCESS.2024.3524415.

\bibitem{ando2026}
Ando, T., Takeuchi, A., Eto, Y., Munakata, Y., \& Nakashima, Y. (2026).
Q-Snap: Quantization-aware dynamic chunking for LLM execution on a CGLA.
\textit{Proceedings of ICISN 2026}.

\bibitem{ando2026_llm}
Ando, T., \& Nakashima, Y. (2026). Energy-efficient Llama3 acceleration
on a CGLA by offloading computational kernels. \textit{IEEE International
System-On-Chip Conference (SOCC2025)} (poster).

\bibitem{ando2025_whisper}
Ando, T., Eto, Y., \& Nakashima, Y. (2025). Energy-efficient hardware
acceleration of Whisper ASR on a CGLA. \textit{CANDAR'25}.
\textbf{Best Paper Award}.

\bibitem{takeuchi2025}
Takeuchi, A., \& Nakashima, Y. (2025). Power efficient attention
acceleration on CGLA. \textit{IEEE International System-On-Chip
Conference (SOCC2025)}.

\bibitem{minamoto2026topology}
Minamoto, M. (2026). On the topology of symbol grounding: Sensory space
structure as a determinant of cognitive representation.
Preprint (not yet on arXiv). Available at:
\texttt{https://github.com/minamominamoto/topology-of-grounding}

\bibitem{harnad1990}
Harnad, S. (1990). The symbol grounding problem. \textit{Physica D:
Nonlinear Phenomena}, 42(1--3), 335--346.

\bibitem{vietoris1922}
Vietoris, L. (1922). Bereiche zweiter Ordnung. \textit{Monatshefte
f\"{u}r Mathematik und Physik}, 32, 258--280.

\bibitem{gromov1981}
Gromov, M. (1981). Structures m\'{e}triques pour les vari\'{e}t\'{e}s
riemanniennes. \textit{Textes Math\'{e}matiques}, 1. CEDIC/Fernand Nathan.

\bibitem{jouppi2017}
Jouppi, N. P., et al. (2017). In-datacenter performance analysis of a
tensor processing unit. \textit{Proceedings of ISCA 2017}, 1--12.

\bibitem{garcez2022}
Garcez, A., \& Lamb, L. C. (2022). Neurosymbolic AI: The 3rd wave.
\textit{Artificial Intelligence Review}, 56, 12387--12406.

\bibitem{bargmann2006}
Bargmann, C. I. (2006). Chemosensation in \textit{C. elegans}.
\textit{WormBook}, 1--29.

\bibitem{white1986}
White, J. G., Southgate, E., Thomson, J. N., \& Brenner, S. (1986).
The structure of the nervous system of the nematode
\textit{Caenorhabditis elegans}. \textit{Philosophical Transactions
of the Royal Society B}, 314(1165), 1--340.

\bibitem{scheffer2020}
Scheffer, L. K., et al. (2020). A connectome and analysis of the adult
\textit{Drosophila} central brain. \textit{eLife}, 9, e57443.

\bibitem{davies2021}
Davies, M., et al. (2021). Advancing neuromorphic computing with Loihi:
A survey of results and outlook. \textit{Proceedings of the IEEE},
109(5), 911--934.

\bibitem{furber2014}
Furber, S. B., Galluppi, F., Temple, S., \& Plana, L. A. (2014).
The SpiNNaker project. \textit{Proceedings of the IEEE}, 102(5), 652--665.

\bibitem{schemmel2010}
Schemmel, J., et al. (2010). A wafer-scale neuromorphic hardware system
for large-scale neural modeling. \textit{Proceedings of ISCAS 2010},
1947--1950.

\bibitem{kriegeskorte2008}
Kriegeskorte, N., Mur, M., \& Bandettini, P. (2008). Representational
similarity analysis---connecting the branches of systems neuroscience.
\textit{Frontiers in Systems Neuroscience}, 2, 4.

\bibitem{biderman2023}
Biderman, S., et al.\ (2023).
Pythia: A suite for analyzing large language models across training and scaling.
\textit{Proceedings of ICML 2023}.

\bibitem{radford2019language}
Radford, A., et al.\ (2019).
Language models are unsupervised multitask learners.
\textit{OpenAI Blog}, 1(8).

\end{thebibliography}

\end{document}
